\subsection{Joint finite-tail matrices}
\label{sec:window-joint-cofinal}

The residual interval of Proposition~\ref{prop:window-inverse-residual}
can be used on a prescribed sequence of growing windows.  We first
write its complete matrix in the unnormalized Legendre coordinates.
We then choose finite indices and rational tail maps for which every
approximation error tends to zero.  The signed low matrix remains the
arithmetic quantity to be estimated.

Fix $T$ and abbreviate $P_n^T$ to $P_n$.  Put $L=2T$, with
$p_n,e_n,g_n$ as in \eqref{eq:window-legendre}.

The following construction applies to a self-adjoint
$\mathcal W$ with the same form domain and with every polynomial
in its operator domain, such that
\[
 \mathcal W\succeq D_T-aI,\qquad
 \|(1-P_n)\mathcal WP_n\|\le b,
 \qquad d_n:=H_n-a>0\quad(n\ge N).
\]
The operator-domain condition is automatic for the bounded
self-adjoint perturbations of $W_T$ used below.
Choose $N<M\le J$, and set $R_{N,M}=P_M-P_N$. Let
$E_N:\mathbb C^N\to L^2$ and
$E_R:\mathbb C^{M-N}\to L^2$ be the synthesis maps for the
unnormalized polynomial bases of $P_N$ and $R_{N,M}$, respectively.
Their Gram matrices are $G_N=\operatorname{diag}(g_n)_{n<N}$
and $G_R=\operatorname{diag}(g_n)_{N\le n<M}$.
Define
\[
 A=E_N^*\mathcal WE_N,\qquad
 B=E_R^*\mathcal WE_N,\qquad
 D=E_R^*\mathcal WE_R.
\]
For any finite matrix $V:\mathbb C^N\to\mathbb C^{M-N}$, put
\begin{align}
 T(V)&=B^*V+V^*B-V^*DV,\\
 r(V)&=B-DV,\\
 U_V&=E_N-E_RV,\\
 Z_n(V)u&=\langle\mathcal W U_Vu,p_n\rangle_2\quad(u\in\C^N),\\
 \Gamma_J(V)&=(\mathcal WU_V)^*(1-P_J)(\mathcal WU_V)\\
 &= (\mathcal WU_V)^*(\mathcal WU_V)
       -\sum_{n=0}^{J-1}\frac{Z_n(V)^*Z_n(V)}{g_n}.
\end{align}
The inner product is linear in its first argument, so $Z_n(V)$
is a linear row.  The last line is a matrix identity on $\C^N$.
It includes every mixed product in the complete operator. Define
\begin{align}
 \rho(V)&=r(V)^*
 \operatorname{diag}\!\left(\frac1{g_nd_n}\right)_{N\le n<M}
 r(V),\\
 \Omega_J(V)&=
 \sum_{n=M}^{J-1}\frac{Z_n(V)^*Z_n(V)}{g_nd_n}
 +\frac{\Gamma_J(V)}{d_J},\\
 F_{N,M,J}(V)&=A-T(V)-\rho(V)-\Omega_J(V).
 \label{eq:cofinal-joint-matrix}
\end{align}

\begin{proposition}\label{prop:cofinal-joint}
If $F_{N,M,J}(V)\succeq-\eta G_N$, with $\eta\ge0$, then
$\mathcal W\succeq-\eta I$ on $\mathcal Q_T$.
\end{proposition}

\begin{proof}
Let $\mathcal K$ be the operator of the restricted tail form
$(1-P_N)\mathcal W(1-P_N)$.  It is strictly positive, and its inverse
is bounded above by the inverse of the diagonal operator with
eigenvalues $d_n$. For $u\in\mathbb C^N$, put
$x=(1-P_N)\mathcal WE_Nu$ and $y=E_RVu$.
The exact inverse identity gives
\[
 \langle x,\mathcal K^{-1}x\rangle
 =2\operatorname{Re}\langle x,y\rangle-\langle y,\mathcal Ky\rangle
 +\langle x-\mathcal Ky,\mathcal K^{-1}(x-\mathcal Ky)\rangle.
\]
The first two terms are $u^*T(V)u$. The retained part of the
residual is $r(V)u$, in covariant polynomial coordinates.
The remote part is
\[
 (1-P_M)\mathcal W(E_N-E_RV)u.
\]
Inverse order bounds the retained part by $u^*\rho(V)u$.
For remote indices, $d_n$ is increasing; Parseval and
$d_n\ge d_J$ for $n\ge J$ give the bound
$u^*\Omega_J(V)u$.
Thus the exact low Schur complement is at least
$F_{N,M,J}(V)$. Completing the square on the closed form domain
gives the assertion, since
$-\eta\|P_Ng\|^2\ge-\eta\|g\|^2$.
\end{proof}

The allowance here is nonpositive.  A positive target still requires
the shifted denominators of \eqref{eq:window-shifts}.

The retained residual $r(V)=B-DV$ and the remote expression
$Z_n(V)=\langle\mathcal WE_N,p_n\rangle_2-
\langle\mathcal WE_R,p_n\rangle_2 V$ are combined before their
Gram matrices are formed. In particular, the products between
these two summands, and between every component of $\mathcal W$,
remain in \eqref{eq:cofinal-joint-matrix}.

For the exact finite solve
\[
 V_0=D^{-1}B,
\]
one has $r(V_0)=0$ and
\begin{equation}\label{eq:cofinal-exact-solve}
 S_M=A-B^*D^{-1}B,\qquad
 F_{N,M,J}(V_0)=S_M-\Omega_J(V_0).
\end{equation}
Here $S_M$ is the exact finite Schur matrix.  The inverse exists
because $D\succeq d_NG_R>0$.
Define the physical norm of a finite map by
\[
 \|V\|_{\mathrm{ph}}
 =\|G_R^{1/2}VG_N^{-1/2}\|.
\]
The joint residual has the uniform bound
\begin{equation}\label{eq:cofinal-joint-residual-rate}
 \|G_N^{-1/2}\Omega_J(V)G_N^{-1/2}\|
 \le \frac{b^2(1+\|V\|_{\mathrm{ph}}^2)}{d_M}.
\end{equation}
Indeed, $(1-P_M)\mathcal WP_M$ has norm at most $b$, and the
orthogonal synthesis $E_N-E_RV$, after normalization of its
source, has squared norm $1+\|V\|_{\mathrm{ph}}^2$.
Every denominator in $\Omega_J$ is at least $d_M$.
This bound is a ceiling on the jointly computed residual; it
does not replace that residual by separate component Grams.

\subsection{Explicit cofinal indices}

Let
\[
 m=j+1\ge2,\qquad T_j=\log m,\qquad B_j=16m^2.
\]
All admitted prime powers satisfy $n<m^2$; a term at $n=m^2$
has zero correlation. Since $\Lambda_{\mathrm{vM}}(n)\le\log n$,
\[
 \nu(T_j)\le4T_j\sum_{n=2}^{m^2}n^{-1/2}
 \le8T_j(m-1)<8m^2.
\]
Using $\pi<4$, $\gamma<1$, $T_j<m$, and
$2\sinh T_j=m-m^{-1}<m$ in \eqref{eq:window-global-comparison} and
\eqref{eq:window-uniform-cross} gives the following bound.
For the scalar logarithm, $\log(2\pi\log m)<\log(8m)<2m$
when $m\ge2$.  Thus each of $\mathcal A(T_j)$ and $b(T_j)$
is less than $8m^2+7m/2+4<16m^2$, so
\begin{equation}\label{eq:cofinal-ab-bound}
 \mathcal A(T_j)<B_j,\qquad b(T_j)<B_j.
\end{equation}
These bounds follow from the local terms of the complete form.

For the exact operator, one possible choice is
\begin{align}
 N_j^{(0)}&=2^{\,4B_j+2},\\
 M_j^{(0)}&=2^{\,4B_j+8B_j^2m^2+2},\\
 J_j^{(0)}&=2M_j^{(0)}.
\end{align}
Since $H_n\ge\log(n+1)$ and $\log2>1/2$,
\[
 H_{N_j^{(0)}}-\mathcal A(T_j)>B_j+1,\qquad
 H_{M_j^{(0)}}-\mathcal A(T_j)>4B_j^2m^2.
\]
Consequently the exact solve satisfies
$\|V_0\|_{\mathrm{ph}}<1$, and
\[
 \|G_N^{-1/2}\Omega_J(V_0)G_N^{-1/2}\|
 <\frac1{2m^2}.
\]
Writing
\[
 \alpha_j^{(0)}=
 \max\{0,-\lambda_{\min}(G_N^{-1/2}S_MG_N^{-1/2})\},
\]
Proposition~\ref{prop:cofinal-joint} gives
\begin{equation}\label{eq:cofinal-exact-bound}
 W_{T_j}\succeq-\left(\alpha_j^{(0)}+\frac1{2m^2}\right)I.
\end{equation}
The finite signed matrix in this statement is exactly
$A-B^*D^{-1}B$, with the complete form in every block.

\subsection{Bounded approximations and finite rational maps}

The same construction accommodates prescribed operator and
coefficient errors. Set
\[
 e_j=\frac1{64m^2}.
\]
Choose a real symmetric approximation $\widetilde W_j$ with
bounded difference $\|W_{T_j}-\widetilde W_j\|\le e_j$.
The construction below supplies this bound and computable polynomial
moments.  Define
\[
 \mathcal W_j=\widetilde W_j-e_jI.
\]
Then
\[
 0\preceq W_{T_j}-\mathcal W_j\preceq2e_jI,\qquad
 \mathcal W_j\succeq D_{T_j}-(B_j+2e_j)I.
\]
Its finite cross blocks have norm at most $B_j+e_j$.
Use the enlarged indices
\begin{align}
 N_j&=2^{\,8B_j+4},\\
 M_j&=2^{\,8B_j+32B_j^2m^2+4},\\
 J_j&=2M_j.
 \label{eq:cofinal-robust-indices}
\end{align}
For $a=B_j+2e_j$ these give
\[
 d_{N_j}>3B_j+1,\qquad
 d_{M_j}>16B_j^2m^2.
\]
The exact finite map for $\mathcal W_j$ therefore has
$\|D^{-1}B\|_{\mathrm{ph}}<1/2$.

A finite polynomial in the retained matrix gives an explicit
approximate solve. In physical coordinates write
\[
 D^\#=G_R^{-1/2}DG_R^{-1/2},\qquad
 B^\#=G_R^{-1/2}BG_N^{-1/2}.
\]
One has $\|D^\#\|\le4M_j$.  To prove this bound, the integral
formula
\[
 P_n(s)=\frac1\pi\int_0^\pi
       \bigl(s+i\sqrt{1-s^2}\cos\theta\bigr)^n\,d\theta
       \qquad(-1\le s\le1)
\]
gives $|P_n(s)|\le1$.  Expanding the generating series of the
integral gives $(1-2sz+z^2)^{-1/2}$ and hence identifies these
polynomials.  Consequently
$\sum_{n<M}|e_n(x)|^2\le M^2/L$, while
\[
 0\le-\tfrac12\log(1-s^2)
 \le-\tfrac12\log(1-|s|)
 \quad\Longrightarrow\quad
 \int_{-T}^TV_T^2\,dx\le L/2.
\]
Summing the squared images of the orthonormal basis gives
\[
 \|V_TP_M\|^2\le\sum_{n<M}\|V_Te_n\|_2^2
       \le M^2/2.
\]
The diagonal part has norm at most $H_{M-1}\le M$.
For $T=\log m$, $|-\gamma-\log(2\pi T)|<B_j$, and
the norms of the other bounded terms sum to less than $B_j$.
It follows that
\[
 \|\mathcal W_jP_M\|
 \le2M+2B_j+2e_j<4M,
\]
using $M\ge2B_j+2$, which holds for the specified indices.
Put
\[
 s_j=\frac1{4M_j},\qquad
 k_j=32mM_j(B_j+1),
\]
and define
\[
 (V_j^{\mathrm{pol}})^\#
 =s_j\sum_{\ell=0}^{k_j-1}(I-s_jD^\#)^\ell B^\#.
\]
The spectrum of $I-s_jD^\#$ lies in
$[0,1-d_N/(4M_j)]$. Hence
\[
 \|B^\#-D^\#(V_j^{\mathrm{pol}})^\#\|
 \le(B_j+1)e^{-k_j/(4M_j)}
 \le\frac1{8m},
 \qquad
 \|(V_j^{\mathrm{pol}})^\#\|<\frac12.
\]
Choose rational coefficients $V_j$ with entrywise error at most
\[
 \frac1{32mM_j^2}
\]
from $V_j^{\mathrm{pol}}$ in the unnormalized polynomial coordinates.
The ratio between the largest retained Gram weight and the
smallest low Gram weight is less than one. Thus the physical
map error is at most $1/(32mM_j)$, giving
\begin{equation}\label{eq:cofinal-finite-forcing}
 \|V_j\|_{\mathrm{ph}}\le1,\qquad
 \|G_R^{-1/2}(B-DV_j)G_N^{-1/2}\|
 \le\frac1{4m}.
\end{equation}
All entries of the finite polynomial map are computable real
numbers. Rational approximants with the displayed error are
obtained from finite enclosures; no choice of an exact rounding
boundary is needed.

Equations \eqref{eq:cofinal-joint-residual-rate} and
\eqref{eq:cofinal-finite-forcing} imply
\begin{equation}\label{eq:cofinal-two-residual-budgets}
 \|G_N^{-1/2}\rho(V_j)G_N^{-1/2}\|
 \le\frac1{16m^2},\qquad
 \|G_N^{-1/2}\Omega_J(V_j)G_N^{-1/2}\|
 <\frac1{4m^2}.
\end{equation}

Let $S_M^W$ and $S_M^{\mathcal W}$ be the exact finite
zero-target Schur matrices for $W_{T_j}$ and $\mathcal W_j$
on these same indices. Minimization over the retained tail gives
\begin{equation}\label{eq:cofinal-schur-perturbation}
 0\preceq
 G_N^{-1/2}(S_M^W-S_M^{\mathcal W})G_N^{-1/2}
 \preceq4e_jI.
\end{equation}
The upper bound follows by inserting the minimizer for
$\mathcal W_j$, whose physical norm is at most one, into the
form for $W_{T_j}$. Moreover,
$A-T(V_j)\succeq S_M^{\mathcal W}$; the difference is the
positive finite inverse residual
$(B-DV_j)^*D^{-1}(B-DV_j)$.
Thus, if the gap between the exact finite matrix and its chosen
lower arithmetic enclosure is at most $1/(16m^2)$ in physical
norm, the construction has total allowance at most
\begin{equation}\label{eq:cofinal-robust-bound}
 \alpha_j+4e_j+\frac1{16m^2}+\frac1{4m^2}+\frac1{16m^2}
 =\alpha_j+\frac7{16m^2},
\end{equation}
where, on the indices \eqref{eq:cofinal-robust-indices},
\[
 \alpha_j=\max\{0,-\lambda_{\min}
       (G_N^{-1/2}S_M^WG_N^{-1/2})\}.
\]

\subsection{Explicit finite coefficient error bounds}

The following error calculation is in physical, orthonormal
coordinates; it translates to the original polynomial coordinates
by the displayed Gram factors. The chosen matrix $V$ is treated as
an exact finite map. Let $v=\|V\|_{\mathrm{ph}}$.
Suppose the full $M$-dimensional compression, the $J$ by $M$
coupling matrix, and the Gram matrix of $\mathcal WP_M$ have
norm errors $\varepsilon_H,\varepsilon_C,\varepsilon_K$.
Their centers determine $\widehat r$, $\widehat Z$, and
$\widehat K_V$, where $Z$ is the jointly formed coupling matrix
for $\mathcal W(E_N-E_RV)$ and $K_V$ its complete Gram matrix.
Put
\[
 \varepsilon_r=(1+v)\varepsilon_H,\qquad
 \varepsilon_Z=\sqrt{1+v^2}\,\varepsilon_C.
\]
The complete joint matrix has an explicit physical norm error
\begin{align}
 \omega\le{}&(1+v)^2\varepsilon_H
 +\frac{2\|\widehat r\|\varepsilon_r+\varepsilon_r^2}{d_N}
 +\frac{(1+v^2)\varepsilon_K}{d_J}\\
 &+\frac{2\|\widehat Z\|\varepsilon_Z+\varepsilon_Z^2}{d_M}
 +\tau_R\|\widehat r\|^2
 +\tau_J\|\widehat K_V\|
 +\tau_Q\|\widehat Z\|^2.
 \label{eq:cofinal-coefficient-error}
\end{align}
Here $\tau_R$ bounds the retained inverse-weight error,
$\tau_J$ the scalar error in $1/d_J$, and $\tau_Q$ the norm
error of the diagonal weights
\[
 w_n=\begin{cases}
 -1/d_J,&n<M,\\
 1/d_n-1/d_J,&M\le n<J.
 \end{cases}
\]
The remote correction is exactly $K_V/d_J+Z^*\operatorname{diag}
(w_n)Z$. Its weight norm is at most $1/d_M$, which proves the
remote part of \eqref{eq:cofinal-coefficient-error}.
The first two terms follow by expanding
$T(V)$ and $r(V)^*\operatorname{diag}(d_n^{-1})r(V)$.
The three input errors include the Gram normalization.  Any
subsequent rounding, including evaluation of the normalized map $V$,
is enclosed separately and added to $\omega$.

An entry radius $r$ in the unnormalized $M$ by $M$ compression
or image Gram gives a physical norm error at most
\[
 \frac{M(2M-1)}{L}r.
\]
For a $J$ by $M$ coupling matrix it gives at most
\[
 \frac{\sqrt{JM(2J-1)(2M-1)}}{L}r.
\]
In particular, for the maps in
\eqref{eq:cofinal-finite-forcing}, at $J=2M$ and
$L=2\log m>1$, the choices
\[
 r_j=\frac1{10^6m^2M_j^4},\qquad
 \tau_R,\tau_J\le\frac1{10^8m^2M_j^4},\qquad
 \tau_Q\le\frac2{10^8m^2M_j^4}
\]
make \eqref{eq:cofinal-coefficient-error} smaller than
$1/(1000m^2)$, hence smaller than the budget in
\eqref{eq:cofinal-robust-bound}.
One may use the bounds
$\|\mathcal WP_M\|\le4M$,
$\|(\mathcal WP_M)^*(\mathcal WP_M)\|\le16M^2$,
$\|\widehat r\|,\|\widehat Z\|\le10M$, and
$\|\widehat K_V\|\le33M^2$ to verify this directly.
These are deliberately coarse coefficient bounds; the actual
joint residual matrices can be much smaller.  Choose subsequent
rounding error below $1/(1000m^2)$ as well.  The total center
error is then below $1/(500m^2)$.

If $\widehat F_j$ is a finite center and
$\|G_N^{-1/2}(F_j-\widehat F_j)G_N^{-1/2}\|\le\omega_j$,
define the certified lower matrix and its signed endpoint by
\[
 \underline F_j=\widehat F_j-\omega_jG_N,\qquad
 \sigma_j=\lambda_{\min}
 (G_N^{-1/2}\underline F_jG_N^{-1/2}).
\]
Then
\[
 \underline F_j\preceq F_j\preceq
 \underline F_j+2\omega_jG_N,\qquad
 W_{T_j}\succeq-\max(0,-\sigma_j)I.
\]
The distinction between a center error and a one-sided enclosure
gap accounts for the factor two. The displayed entry radii give
$2\omega_j<1/(16m^2)$, as required in
\eqref{eq:cofinal-robust-bound}. In particular,
\[
 \max(0,-\sigma_j)\le\alpha_j+\frac7{16m^2}.
\]
The endpoint $\sigma_j$ is computed from the single joint matrix
\eqref{eq:cofinal-joint-matrix}, not from separate arithmetic,
polar, or regular-kernel lower bounds.
If an additional scalar eigenvalue enclosure is used, its error
is included in the same one-sided arithmetic budget. Alternatively
one verifies the desired Loewner lower bound directly.

If a finite matrix is instead enclosed in ordinary coefficient
norm by $E_FI$, its contribution to a physical allowance is
\[
 \frac{E_F}{g_{\min}}=\frac{(2N-1)E_F}{L}.
\]
This factor cannot be omitted on a growing sequence of indices.
Dividing all moment matrices and all Gram matrices by the common
factor $L$ gives the equivalent unitary-coordinate formulation
with Gram weights $(2n+1)^{-1}$.

\subsection{Existence of the prescribed bounded approximations}

The error $e_j$ can be achieved without changing the form domain
or any prime-power term. The function $R$ has a global Lipschitz
bound $|R'|\le16$. On $[0,1]$, the coefficient estimate from
\eqref{eq:window-r-error} gives
\[
 |R'(r)|\le17\sum_{n\ge1}n3^{-n}=51/4.
\]
On $[1,\infty)$, use
$J(r)=e^{-r/2}/(1-e^{-2r})$ and $e^{-2}<1/7$.  Differentiation
gives $|J'|\le7/12+7/18$.  Hence
$|R'|\le7/12+7/18+1/2<2$ there.
Use a constant approximation to $R$ on intervals of length
\[
 h_j=\frac{e_j}{96m},
\]
with its value at each midpoint enclosed to error at most
$e_j/(12m)$. Since $L<2m$, its integral-operator error is at
most $e_j/3$. There are at most $192m^2/e_j+1$ intervals.
The resulting real symmetric kernel is bounded.

For each polar exponential, take the Taylor polynomial through
degree $k-1$, where
\[
 k=3m+\left\lceil\log_2(48m^2/e_j)\right\rceil.
\]
The factorial bound $k!\ge(k/e)^k$, $e<3$, and $T_j<m$ give
uniform exponential error at most $\sqrt m\,2^{-k}$.
The cross-polar operator error is at most
$16m^2\,2^{-k}\le e_j/3$. The two cross rank-one pieces are
retained.

For Euler's constant, the trapezoid error for $x^{-1}$ gives
\begin{equation}
 H_q-\log q-\frac1{2q}\le\gamma\le
 H_q-\log q-\frac1{2q}+\frac1{8q^2}.
 \label{eq:cofinal-gamma-trapezoid}
\end{equation}
Indeed the error on $[k,k+1]$ is nonnegative and at most
$\frac18\int_k^{k+1}2x^{-3}\,dx$.  Sum these inequalities and
pass to infinity.  At $q=N_\gamma=16m$, the displayed interval
has width $e_j/32$.
Enclosing $\log N_\gamma$ to width at most $e_j/32$ and taking
the rational midpoint gives a scalar gamma error less than
$e_j/3$. The three errors sum to at most $e_j$.
The logarithmic wall and every prime-power shift remain exact.

For every fixed $j$, these approximations send a polynomial to
a finite piecewise polynomial plus its logarithmic-wall product.
The complete image Grams therefore reduce to polynomial moments,
single-logarithm moments on finitely many intervals, and global
squared-logarithm moments. All endpoints and coefficients are
computable, and finite rational enclosures with the prescribed
entry radii exist. This assertion supplies the error budgets;
it supplies no sign for the resulting finite matrix.
The precision thresholds are mathematical requirements on the
finite enclosures. They do not bound the arithmetic complexity of
forming them. The finite polynomial construction of $V_j$ may
use finer intermediate enclosures before its final coefficient
rounding; its stated forcing bound remains the verification
condition for that map.

\subsection{The remaining signed inequality}

Let $\mu_M$ be the lowest variational eigenvalue of the complete
compression $P_MW_TP_M$, and let $\alpha$ be the negative part
of its finite Schur matrix as above. If the retained tail is
positive, then
\begin{equation}\label{eq:cofinal-alpha-variational}
 (-\mu_M)_+\le\alpha
 \le(1+\|D^{-1}B\|_{\mathrm{ph}}^2)(-\mu_M)_+.
\end{equation}
The first inequality follows by completing the square. For the
second, insert the minimizing vector
$E_Nu-E_RD^{-1}Bu$ in the finite variational inequality.
On either explicit family of indices, the multiplicative factor
is at most two.

\begin{theorem}[Equivalent signed estimates on an explicit family]
\label{thm:cofinal-equivalence}
Use the windows, indices, rational maps and enclosures constructed
above.  Let $\alpha_j$ be the negative part of the exact finite
Schur matrix for $W_{T_j}$, measured in the physical norm.
Let $\sigma_j$ be the lowest eigenvalue of the enclosed joint
matrix.  The following conditions are equivalent:
\begin{enumerate}
\item $L_\zeta(f^*\star f)\ge0$ for every $f\in\cA_\zeta$.
\item $\alpha_j\to0$.
\item $(-\mu_{M_j}(T_j))_+\to0$.
\item $\max(0,-\sigma_j)\to0$.
\end{enumerate}
For every $j$, without assuming any of these conditions,
\begin{equation}
 Q_{T_j}(g)\ge-\max(0,-\sigma_j)\|g\|_2^2,
 \qquad
 \max(0,-\sigma_j)\le\alpha_j+\frac7{16(j+1)^2}.
 \label{eq:cofinal-final-allowance}
\end{equation}
\end{theorem}
\begin{proof}
The two inequalities are the enclosed joint bound and
\eqref{eq:cofinal-robust-bound}.  Inequality
\eqref{eq:cofinal-alpha-variational} proves the equivalence
of the second and third conditions.  The second implies the
fourth by \eqref{eq:cofinal-final-allowance}.
If the fourth holds, extend any fixed $g\in\mathcal Q_T$ by
zero to the cofinal windows.  Its form value and norm remain
unchanged by Proposition~\ref{prop:window-extension}.
Taking the limit gives $Q_T(g)\ge0$.  The cutoff argument in
Theorem~\ref{thm:window-cofinal} then proves the first condition.
Conversely, positivity on $\cA_\zeta$ gives positivity on
compact smooth tests and hence on $\mathcal Q_T$ by the form-core
property.  Every exact finite Schur matrix is then nonnegative,
so $\alpha_j=0$ for every $j$.
\end{proof}

All operator, finite-solve, tail and coefficient errors in this
construction have explicit vanishing bounds.  The remaining
signed condition is equivalent to full Weil positivity.
Compactness and a positive tail do not establish it.

For each fixed $T$, polynomial form-core density gives
$\mu_M(T)\downarrow\inf\operatorname{spec}W_T$. This fixed-window
convergence is distinct from a signed estimate as $T\to\infty$.
The explicit indices above give uniform absolute control of the
discarded corrections while leaving the exact signed low
inequality visible. No positive margin in $T$ is asserted.
In particular, the near-null families in Proposition~\ref{prop:window-near-null}
exclude a positive power-law margin without contradicting these
vanishing negative allowances.
